\documentclass[11pt]{article}


\addtolength{\textwidth}{2cm}
\addtolength{\oddsidemargin}{-1cm}

\addtolength{\textheight}{0.7cm}
\addtolength{\topmargin}{-0.35cm}


\parskip 0.15cm





%%%%%%%%%%%%%%%%%%%%%%%%%%%%%%%%%%%%%%%%%%%%%%%%%%%%%%%
% definitions specific to this document


\newcommand{\heading}{\Large}  

%%%%%%%%%%%%%%%%%%%%%%%%%%%%%%%%%%%%%%%%%%%%%%%%%%%%%%%


\begin{document}



\begin{center}
{\sc\LARGE Research Statement}
\end{center}

\vspace{0.5cm}



My PhD training has been in condensed matter theory.  I am interested
in studying quantum many-body phenomena in solids, trapped atomic
gases, artificial structures like quantum wells and dots, and other
systems that display quantum coherence or collective phenomena.

My research during my PhD has been in three main areas.  Work with my
advisor, Andrei E Ruckenstein, has focused on \emph{Bose-Einstein
condensation} (BEC) phenomena. In collaboration with graduate student
colleagues, I studied \emph{Wigner crystal physics}.  My latest
project is a study of \emph{exciton condensation in coupled quantum
wells}.  Although I will be pursuing several loose collaborations,
this last project is primarily my own.


{\heading BEC Transition Temperature}

My first major project, with Andrei Ruckenstein, involved the
well-known ``Tc problem''.  Calculation of the BEC transition
temperature of a weakly repulsive Bose gas has continued to elude
theorists, forty years after the first attempts.  Our approach was to
use second-order perturbation theory in the two-body t-matrix.

Using a quasiparticle description of the un-condensed Bose gas, we
reduced the calculation of Tc to a calculation of the quasiparticle
spectrum at the critical point.  Our main contribution is a method of
regularizing the infrared divergences appearing in the description of
the transition point, resulting in a shift of Tc approximately
proportional to a*sqrt(-ln[a]), where a is the scattering length.
Our analysis points strongly to a non-analytic contribution in the
dependence of Tc on the interaction.

The main results have been submitted, cond-mat/0212590, and a
collection of shorter calculations will be posted as a comment in
January.


{\heading Variational Description of Bose Condensate}.

For the past several months, my work with my advisor has concentrated
on a variational description of the condensed Bose gas.  The
variational wavefunction used is a squeezed coherent state, using
ideas borrowed from Quantum Optics. Development of an alternate
description is expected to provide additional insight into interaction
effects in the weakly non-ideal Bose gas.

In developing this formalism, I calculated the variational parameters
and the excitation spectrum.  More importantly, I worked out the
details of the manifestation of U(1) gauge invariance and
gaplessness (Goldstone theorem), which are essential to any
description of the condensed state.

Two directions are being explored.  The first is an attempt to add
more freedom to the variational function in a way that will
incorporate effects beyond mean field.  Secondly, the formalism is
being applied to a description of interaction effects in matter-wave
interference phenomena, in particular four-wave mixing of Bose
condensates.


{\heading Wigner Crystal Physics}.

In addition to the work done with my advisor, in spring 2001 I took
the initiative to begin work on an independent research project, in
collaboration with a graduate student colleague (S. Pankov).  We
started studying the AC response of a disorder-pinned Wigner crystal,
under high magnetic fields.  This project eventually turned into a
consideration of Wigner crystal structures on a cryogenic liquid
(e.g., helium).  I later invited another PhD student (I. Paul) to join
the collaboration.


For a 2D electron system on a liquid surface, surface distortions
introduce additional physics, and can be tuned by varying the pressing
electric field ($E_\perp$) used to hold the electrons to the surface.
We found that, at densities low enough such that the inter-electron
distance is comparable to the capillary length, there is at least one
structural transition, so that a lattice geometry other than the usual
triangular one becomes more stable.  The transition is driven by the
competition between two central forces between the electrons, the
long-range Coulomb repulsion and a short-range attraction, the latter
being due to surface distortions.  This result has been submitted for
publication (cond-mat/0211648).

At present, our work with this system is evolving in two directions.

One program is a continuation of our previous work --- to map out
exactly which structures are the stable ones in the various regions of
the $E_\perp$-density phase diagram.  In addition, the system has a
crossover from the single-electron lattices at low electron densities,
to a charge-density-wave structure at higher densities.  The details
of this crossover are entirely unknown.  Our current approach (myself
and I. Paul) to both issues is to study the dynamical matrix for a
number of lattice geometries.

I am particularly enthusiastic about the second direction, which is
inspired by experimental attempts to achieve quantum melting of a
Wigner crystal on a thin helium film.  Because of the attraction due
to exchange of ripplons (quantized surface vibrations), the electron
liquid in the quantum regime is expected to have a Cooper instability.
With S. Pankov, I am examining the BCS transition for 2D electrons on
thin helium films.



{\heading  Excitons in Biased Double Quantum Well}.

My latest project, on which I plan to report during the 2003 APS March
meeting, involves Bose condensation of excitonic systems in coupled
quantum wells in semiconductors.  This is essentially an independent
project, and I intend to devote a significant part of the next year
working on it.  

The idea of producing an excitonic gas, dense enough and cold enough
to condense, has experienced its latest revival in the context of
double quantum wells.  The wells are biased with an electric field so
that the electrons and holes are spatially separated, leading to
excitons with a longer lifetime than 3D excitons, and a faster cooling
rate.  Intriguing features of this system include the effect of laser
light (used to create the exciton gas) on the condensate, the BCS -
BEC crossover in the description of the condensation process, and the
inherently non-equilibrium nature of the system.  None of these
aspects have been worked out in detail in the context of the 2D
coupled-well system.

My broad aim is (I) to adapt calculations done for 3D and single-well
systems to the double-well case, and (II) to investigate phenomena
specific to the 2D coupled-well case (details of Kosterlitz-Thouless
transition, condensation in 2D traps).  As a first calculation, I am
examining the order parameter enhancement due to the presence of a
coherent (classical) laser field.

Collaboration with Piers Coleman is focusing on non-equilibrium
aspects.  I have also started discussions with Morrel H. Cohen and
Phil Platzman (theorists), and am scheduling a visit to David Snoke's
lab in Pittsburgh, where an experimental effort for Bose condensation
of excitons in coupled wells is under way.




{\heading Future Directions}.

In addition to the ongoing work, there are several ideas that I intend
to follow up on, during the next year or so.
%
\begin{enumerate}
%
\item For the inter-leaved vortex lattices in a rapidly rotating
two-component BEC, several structural transitions have been predicted,
using a Quantum Hall formalism.  It would be interesting to
re-interpret these transitions in terms of an inter-vortex interaction
potential.  This would also allow me to use calculations and code from
my Wigner crystal work to study these transitions.
%
\item 2D electrons on helium --- 
%
I have been communicating with the Rutgers experimental group
(I. Skachko \& E.Y. Andrei) that studies 2D electrons on \emph{thin}
films, and I would like to pursue a closer collaboration.  I am
interested in suggesting experiments probing the Wigner crystal in the
presence of a magnetic field.  My personal motivation is to learn more
about the physics of 2D electrons in a magnetic field.  The additional
effects due to tiny substrate thickness make this project more
attractive.
%
\item BEC interference ---
%
One aspect common to several classes of BEC experiments is the
presence of two or more macroscopically occupied momentum states,
leading to matter-wave interference phenomena.  One approach is a
Bogoliubov-type calculation for the excitation spectrum in the
presence of two condensates.  I have some preliminary results, which I
would like to explore further.
%
\item Other approaches to the Tc problem --- 
\begin{enumerate}
%
\item One plan is to use finite-size scaling ideas, in conjunction with
second-order perturbation theory in the grand-canonical ensemble, to
calculate the transition temperature Tc for the weakly non-ideal
Bose gas.
%
\item A more ambitious aim is to calculate Tc using perturbation
theory in the full many-body T-matrix (sum of ladder diagrams).
This approach is expected to vastly improve the quasiparticle
description used previously.  My first attempt, near the beginning of
my PhD, was unsuccessful.  I would like to attack the problem again,
now being much more familiar with many-body techniques.
%
\end{enumerate}
%
%\item  Atoms in optical lattices --- 
%
\item Light scattering off BECs --- Earlier, I studied some details of
the light scattering experiments off atomic Bose condensates.  Now
that I am learning Keldysh (non-equilibrium) techniques in the context
of photonic effects on excitonic condensates, I would like to
re-examine the light scattering experiments using this formalism, to
calculate quantities such as non-equilibrium distribution functions
and super-radiance thresholds.
%
\end{enumerate}



\vspace{0.2cm}

In summary, I plan to continue working on quantum many-body phenomena,
in systems where inter-particle correlations and interactions have
non-trivial effects.  Apart from continuing projects started at
Rutgers, I am eager to start work on new systems and models, and I
intend to actively pursue and organize fresh collaborations.


Over the last few years, there has been considerable progress in
creating artificial nanoscale structures, and in performing novel
experiments with laser-cooled atomic systems.  The synthesis of new
materials continues to produce new surprises and challenges to
theorists.  Condensed-matter theory, the science of emergent phenomena
in many-particle systems, is as vibrant as ever.  I have enjoyed
immensely my studies of this fascinating field, and wish to continue
and broaden my investigations of correlated matter.


\vspace{2cm}

\begin{flushright}
Masudul Haque\\
December 2002
\end{flushright}


\end{document}

